\documentclass[12pt,a4paper]{article}

\usepackage{fontspec}
\usepackage{xcolor}
\usepackage{geometry}
\usepackage{makeidx}
\usepackage{biblatex}

\geometry{
	margin=1in
}

\setmainfont{DejaVu Serif}
\setsansfont{DejaVu Sans}
\setmonofont{Courier New}

\definecolor{pagebg}{RGB}{255,242,230}
\pagecolor{pagebg}

\addbibresource{references.bib}

\makeindex

\renewcommand{\thesection}{\Alph{section}}

\newsavebox{\messagebox}

\begin{document}
	
	\begin{center}
		{\Large\textbf{Assignment 6 - Document Structure}}
		
		\vspace{0.3cm}
		
		TW Lab 2025
	\end{center}
	
	\vspace{1cm}
	
	\begin{center}
		Write \LaTeX{} code to output the following sections.
		Also follow any instructions provided in the document.
		Use \LaTeX{} reference and appropriate packages
		wherever required.
	\end{center}
	
	
	\section{Fonts and background}
	
	Try to match the page background color. Then create the
	following trivlist:
	
	\begin{trivlist}
		
		\item This text is written in DejaVu Serif, the main font.
		
		\item {\sffamily This text is written in DejaVu Sans,
			the sans-serif font.}
		
		\item {\ttfamily This text is written in Courier,
			the monospace font.}
		
	\end{trivlist}
	
	Use the package \texttt{fontspec} and make sure that you
	configure TeXstudio to use XeLaTeX as the default compiler.
	
	
	\section{Spacing}
	
	This is the first line of a paragraph.\\
	This is another line, but not another paragraph.
	
	\vspace{1cm}
	
	This paragraph starts 1cm below. Now, observe both the
	paragraphs given below and how they differ in hyphenation:
	
	\begin{quote}
		$>$ This is another paragraph. Incomprehensibilities
		and antidis-establishmentarianisms are two long words.
	\end{quote}
	
	\begin{quote}
		$>$ This is another paragraph.
		Incomprehensibilities and antidisestablishmentarianisms
		are two long words.
	\end{quote}
	
	Recreate the paragraphs above and then continue from
	next page, leaving space for section C.
	
	\index{background}
	
	\stepcounter{section}
	
	\newpage
	
	
	\section{Numbering}
	
	Ensure that all sections are numbered using alphabets.
	
	\[
	x^2+y^2=1
	\tag{1}
	\]
	
	Equation \ref{eq:circle} denotes a circle. Now add the
	equation of an ellipse before the equation of the circle.
	
	\begin{equation}
		\frac{x^2}{a^2}+\frac{y^2}{b^2}=1
		\label{eq:ellipse}
	\end{equation}
	
	\begin{equation}
		x^2+y^2=1
		\label{eq:circle}
	\end{equation}
	
	Equation \ref{eq:ellipse} denotes an ellipse and
	Equation \ref{eq:circle} denotes a circle.
	
	\index{numbering}
	
	
	\section{Boxes and Lines}
	
	\savebox{\messagebox}{
		This message is written in a savebox.
	}
	
	\usebox{\messagebox}
	
	\usebox{\messagebox}
	
	\usebox{\messagebox}
	
	\vspace{0.5cm}
	
	\parbox{5cm}{
		This message is wrapped in a 5cm paragraph box.
	}
	
	\vspace{0.3cm}
	
	\noindent
	\rule{0.5em}{1em}
	
	\vspace{0.1cm}
	
	\noindent
	\fbox{
		Note that these boxes are used for highlighting
		and layout design.
	}
	
	\index{savebox}
	
	\vspace{0.5cm}
	
	Create an index of some relevant terms from this document.
	Also add the bibliography [2] as shown in the following
	pages. Add the missing entry to the Latex-in-short book
	[1], in the bib file attached and cite it appropriately.
	Carefully understand how index file and bib files are
	compiled.
	
	I used the short introduction to \LaTeX{} as a reference
	while preparing this document \cite{lshort}. Information
	about bibliography management was also referred to
	\cite{overleaf}.
	
	\clearpage
	
	\printindex
	
	\clearpage
	
	\printbibliography[title={References}]
	
\end{document}